\chapter*{Conclusion générale}
\addcontentsline{toc}{chapter}{Conclusion générale}
\label{chap:conclusion}

% ------------------------------------------------------------------
\section*{Rappel de la démarche}

Ce mémoire s'est donné pour objectif d'évaluer, par simulation, la conformité et la
robustesse du module RFM de la plateforme SaaS ELMES-TEKA, conçue pour répondre aux besoins
de gestion commerciale des petits commerçants congolais formels et informels. L'artefact
évalué — la fonction \texttt{computeRFM} au commit
\texttt{b062a93a7be96d082dee9afe66b9f4f1f9ebecaa} — classe chaque client dans l'une des cinq
catégories RFM par application d'une règle de tercilisation sur les dimensions Récence,
Fréquence et Montant.

La démarche adoptée relève de la \emph{Design Science Research} \citep{Hevner2004,
Peffers2008}~: ELMES-TEKA est un artefact construit pour résoudre un problème réel, et la
présente campagne constitue l'étape \emph{evaluate} du cycle de vie de cet artefact. Le
raisonnement hypothético-déductif a guidé la formulation de deux hypothèses préalables à
l'analyse, dont les critères de décision ont été fixés a priori. La simulation Monte-Carlo,
structurée selon le cadre ADEMP \citep{Morris2019}, a été conduite sur 12\,000 essais
produits par croisement de quatre tailles de corpus ($n \in \{3, 10, 50, 200\}$), trois
profils de distribution (\texttt{regular}, \texttt{skewed}, \texttt{ties}) et dix scénarios
de perturbation, avec 100 répétitions par cellule et graines explicitement consignées pour
assurer la reproductibilité complète des résultats.

% ------------------------------------------------------------------
\section*{Réponse à l'hypothèse H1 — Conformité}

La première hypothèse portait sur la conformité du moteur TypeScript à une implémentation
de référence en Python calculant les sommes monétaires en centimes entiers. Sur les
12\,000 essais de la campagne, 200 présentent au moins un écart~: le taux de non-conformité
est de $1{,}667\,\%$ (intervalle de confiance Monte-Carlo à $95\,\%$ : $[1{,}448\,\% ;\,
1{,}904\,\%]$). Ce résultat est inférieur au seuil de tolérance de $2\,\%$ retenu pour le
contexte de PME informelles congolaises, dont les données d'entrée sont elles-mêmes soumises
à des imprécisions de saisie~: la décision retenue pour H1 est donc \textbf{conforme sous
réserve}.

La cause des non-conformités est identifiée avec précision~: l'addition des montants de
commandes est effectuée en virgule flottante IEEE\,754 double précision dans le moteur
TypeScript, alors que la référence Python additionne les mêmes montants convertis en centimes
entiers \citep{Haissinsky2012, Ribatet_nd}. Lorsque plusieurs commandes sont agrégées,
l'erreur d'arrondi cumulée peut modifier la relation d'ordre entre deux clients dont les
sommes affichées sont identiques, déplaçant ainsi le score $M$ et, dans deux cas sur 12\,000
essais, l'étiquette de segment. Ce mécanisme est systématique et reproductible~; il ne
relève pas du bruit aléatoire mais d'une fragilité structurelle corrigible.

La mention «~conforme sous réserve~» traduit la prudence épistémologique qui s'impose ici~:
la conformité est établie dans les limites du périmètre testé et sous la cause identifiée.
Elle ne pourra être portée à «~conforme sans réserve~» qu'après application d'un correctif
arithmétique — arrondi au centime avant comparaison de quantiles, ou adoption d'une
bibliothèque de calcul monétaire en entiers telle que \texttt{dinero.js} — suivi d'une
réévaluation ciblée sur les essais ayant produit des écarts. Cette conformité algorithmique
n'implique par ailleurs pas la pertinence prédictive des segments~: la question de savoir
si la tercilisation RFM constitue le meilleur modèle pour anticiper les comportements de
réachat des clients congolais \citep{Fader2005} reste entière et relève d'une validation
de contenu distincte \citep{Salminen2023}.

% ------------------------------------------------------------------
\section*{Réponse à l'hypothèse H2 — Robustesse}

La seconde hypothèse portait sur la robustesse des affectations de segment face à un bruit
monétaire de $\pm1\,\%$, pour les clients suffisamment éloignés des six seuils de tercile
initiaux selon une règle à $5\,\%$. Les quatre cellules admissibles — tailles
$n \in \{50, 200\}$, profils \texttt{regular} et \texttt{skewed} — présentent toutes un
taux de changement moyen de $0{,}000\,\%$ avec des intervalles de confiance de largeur
nulle $[0\,\% ;\, 0\,\%]$. La borne supérieure de chaque intervalle est strictement
inférieure au seuil de $5\,\%$ fixé par le protocole~: la décision retenue pour H2 est
\textbf{compatible}.

Ce résultat s'explique par la construction même du critère d'éloignement~: un client dont
la distance normalisée à chacun des six seuils initiaux est supérieure à $5\,\%$ dispose
d'une marge de sécurité qui absorbe mécaniquement un bruit monétaire de $\pm1\,\%$ sans
franchir le seuil le plus proche. La campagne confirme empiriquement cette garantie
théorique sur 100 répétitions par cellule, sans jamais l'observer mise en défaut dans les
conditions testées.

Deux nuances s'imposent. D'une part, l'intervalle $[0\,\% ;\, 0\,\%]$ traduit une absence
de variation observée et non une démonstration formelle qu'aucun changement ne peut se
produire dans l'ensemble des configurations d'usage possibles~; la robustesse mesurée est
conditionnelle au niveau de bruit ($\pm1\,\%$), aux tailles de corpus ($n \geq 50$) et aux
profils non-ties. D'autre part, le profil \texttt{ties} — tous les clients identiques sur
les trois dimensions RFM — constitue un cas hors domaine exclu du critère H2 par décision
protocolaire antérieure à l'analyse~; sous perturbation de dates, ce profil produit des
taux de changement de 42 à 63\,\% qui reflètent l'indétermination mathématique des
quantiles sur des valeurs identiques, et non un défaut de l'algorithme.

% ------------------------------------------------------------------
\section*{Limites de l'évaluation}

Quatre limites principales bornent la portée des résultats présentés dans ce mémoire.

Premièrement, les corpus utilisés dans la simulation sont synthétiques et non calibrés sur
des données de transactions réelles de petits commerçants congolais. Les distributions de
montants, les fréquences d'achat et les horizons temporels générés peuvent s'éloigner
sensiblement des distributions observées en production~; les résultats décrivent le
comportement de la fonction sur ces mécanismes générateurs spécifiques et ne constituent
pas une évaluation de son comportement en conditions réelles.

Deuxièmement, la référence Python n'a pas été certifiée par une tierce partie indépendante.
Elle a été construite pour les besoins de la campagne à partir de la spécification de
\texttt{computeRFM} et vérifiée sur neuf cas manuels couvrant les configurations limites
critiques~; cette vérification ne couvre cependant pas l'ensemble de l'espace des entrées.
Une erreur propre de la référence non détectée par comparaison avec le moteur TypeScript
ne peut être exclue.

Troisièmement, l'évaluation porte sur la fonction \texttt{computeRFM} de façon isolée,
hors des couches MongoDB, des routes HTTP, du système d'authentification multi-tenant et
des mécanismes de conversion de devises d'ELMES-TEKA. Les effets économiques de la
segmentation sur les décisions commerciales des commerçants n'ont pas été évalués~; ces
dimensions font partie de l'artefact complet mais dépassent le périmètre défini par H1
et H2.

Quatrièmement, les résultats sont valables pour la logique de segmentation telle que codée
au commit \texttt{b062a93a}~: toute modification du code — changement du nombre de groupes,
remplacement des terciles par d'autres quantiles, ajout d'une couche d'arrondi monétaire —
invaliderait les mesures obtenues et nécessiterait une nouvelle campagne de validation.

% ------------------------------------------------------------------
\section*{Perspectives}

Plusieurs directions de travail permettraient d'étendre et de consolider la présente
évaluation.

\paragraph{Validation sur données réelles.}
La priorité la plus directe est la validation du module sur des données de transactions
réelles anonymisées, extraites de la base de production d'ELMES-TEKA. Une telle validation
permettrait de calibrer les mécanismes générateurs de la simulation sur les distributions
effectives et de mesurer la sensibilité du module dans les conditions d'usage réel
\citep{Salminen2023}. Elle soulèverait toutefois des contraintes de confidentialité liées
à la nature multi-tenant de la plateforme, qui devront être traitées avec soin.

\paragraph{Extension multi-devises.}
La cohabitation du dollar américain et du franc congolais dans le commerce kinois constitue
une priorité métier. Le comportement hors domaine observé lors du mélange de devises
(montants additionnés sans séparation~: $10 + 28\,000 = 28\,010$) montre que la fonction
traite actuellement toutes les valeurs comme commensurables, produisant des montants agrégés
économiquement incohérents. L'extension multi-devises nécessiterait de définir une unité de
compte commune et d'intégrer la conversion dans le pipeline de calcul, avec ses implications
sur la robustesse de la segmentation face aux fluctuations de taux.

\paragraph{Validation prospective des segments.}
Une validation prospective comparerait les segments produits par \texttt{computeRFM} aux
comportements de réachat ultérieurs des clients, testant ainsi la validité prédictive des
cinq étiquettes \citep{Fader2005}. Ce type de validation excède la portée d'une évaluation
de conformité algorithmique~; il est néanmoins indispensable pour établir la valeur
décisionnelle des segments dans le contexte opérationnel. Des modèles stochastiques de
valeur client tels que Pareto/NBD \citep{Fader2005} pourraient fournir une référence externe
pour juger de la qualité discriminante des catégories retenues.

\paragraph{Analyses de sensibilité globale et stabilité temporelle.}
L'incorporation d'analyses de sensibilité globale \citep{Lakens2013} permettrait d'estimer
la contribution relative des dimensions $R$, $F$ et $M$ à la variabilité inter-sessions de
la segmentation, au-delà de l'analyse par scénario conduite dans la présente campagne. De
même, une étude de la stabilité temporelle des segments — mesure dans laquelle un client
conserve la même étiquette d'une période à l'autre en l'absence de changement réel de
comportement — constituerait un indicateur pratique directement exploitable par les
commerçants \citep{Mingmongkhol2023}.

\paragraph{Correction arithmétique et réévaluation.}
À court terme, la correction du mécanisme de virgule flottante identifié par H1 devrait
être suivie d'une réévaluation ciblée sur les 200 essais ayant produit des écarts, afin de
confirmer que le taux de non-conformité est bien ramené à zéro et que la décision peut être
portée à «~conforme sans réserve~».

\bigskip

En définitive, ce mémoire établit que le module RFM d'ELMES-TEKA est, dans les conditions
testées, conforme à sa spécification dans $98{,}3\,\%$ des cas et robuste aux perturbations
monétaires de $\pm1\,\%$ pour les clients éloignés des seuils de tercile. Ces deux résultats,
obtenus par une démarche reproductible et traçable, constituent une base solide pour la
prise de décision des développeurs et des responsables produit d'ELMES-TEKA, ainsi qu'un
point de départ documenté pour les travaux de validation ultérieurs sur données réelles.
